\documentclass{amsart}
\usepackage{amsmath,amssymb,amsthm,hyperref}
\newtheorem{theorem}{Theorem}
\newtheorem{lemma}{Lemma}
\newtheorem{proposition}{Proposition}
\theoremstyle{definition}
\newtheorem{definition}{Definition}
\begin{document}

\title{No degree-$4$ Sum-of-Squares refutation rules out seven MUBs in $\mathbb{C}^6$}
\maketitle

\begin{abstract}
We answer the question in the negative: there is \emph{no} degree-$4$ Sum-of-Squares
(SoS) refutation of the polynomial system encoding the existence of seven mutually
unbiased bases (MUBs) in $\mathbb{C}^6$. The proof exhibits an explicit degree-$4$
\emph{pseudo-expectation}, realized as the genuine expectation over seven independent
Haar-random orthonormal bases of $\mathbb{C}^6$. By SoS--moment duality, the existence of
such a functional precludes any degree-$4$ refutation.
\end{abstract}

\section{The polynomial system and the refutation form}\label{sec:setup}

\subsection{Variables and constraints}
We work with $7$ ordered orthonormal bases of $\mathbb{C}^6$. For $k\in\{1,\dots,7\}$
(the basis index), $i\in\{1,\dots,6\}$ (the vector index within a basis) and
$a\in\{1,\dots,6\}$ (the coordinate index in $\mathbb{C}^6$), write the complex entry
\[
v^{(k)}_{i,a}=x^{(k)}_{i,a}+\sqrt{-1}\,y^{(k)}_{i,a},
\qquad x^{(k)}_{i,a},\,y^{(k)}_{i,a}\in\mathbb{R}.
\]
Thus the system has $7\cdot 6\cdot 6\cdot 2 = 504$ real indeterminates, which we collect
into a single real vector
\[
X=\bigl(x^{(k)}_{i,a},\,y^{(k)}_{i,a}\bigr)_{k,i,a}\in\mathbb{R}^{504}.
\]
We write $\langle u,w\rangle=\sum_{a=1}^{6}\overline{u_a}\,w_a$ for the Hermitian inner
product on $\mathbb{C}^6$, conjugate-linear in the first argument. The MUB system consists
of the following polynomial equalities in $X$.

\medskip
\noindent\textbf{(i) Orthonormality (degree $2$).} For each $k$ and each $i$,
\begin{equation}\label{eq:norm}
n^{(k)}_{i}(X)\;:=\;\sum_{a=1}^{6}\Bigl(\bigl(x^{(k)}_{i,a}\bigr)^2+\bigl(y^{(k)}_{i,a}\bigr)^2\Bigr)-1\;=\;0,
\end{equation}
and for each $k$ and each $i\ne j$, splitting $\langle v^{(k)}_i,v^{(k)}_j\rangle$ into real
and imaginary parts,
\begin{equation}\label{eq:orth}
o^{(k),\mathrm{Re}}_{i,j}(X):=\operatorname{Re}\langle v^{(k)}_i,v^{(k)}_j\rangle=0,
\qquad
o^{(k),\mathrm{Im}}_{i,j}(X):=\operatorname{Im}\langle v^{(k)}_i,v^{(k)}_j\rangle=0 .
\end{equation}
Each of $n^{(k)}_i,o^{(k),\mathrm{Re}}_{i,j},o^{(k),\mathrm{Im}}_{i,j}$ is a real polynomial
of degree exactly $2$. Equations \eqref{eq:norm}--\eqref{eq:orth} say precisely that, for
every fixed $k$, the vectors $v^{(k)}_1,\dots,v^{(k)}_6$ form an orthonormal basis of
$\mathbb{C}^6$.

\medskip
\noindent\textbf{(ii) Unbiasedness (degree $4$).} For all $k\ne\ell$ and all $i,j$,
\begin{equation}\label{eq:unbias}
u^{(k,\ell)}_{i,j}(X)\;:=\;\bigl|\langle v^{(k)}_i,v^{(\ell)}_j\rangle\bigr|^2-\tfrac16\;=\;0 .
\end{equation}
Since $\bigl|\langle v^{(k)}_i,v^{(\ell)}_j\rangle\bigr|^2
=\bigl(\operatorname{Re}\langle v^{(k)}_i,v^{(\ell)}_j\rangle\bigr)^2
+\bigl(\operatorname{Im}\langle v^{(k)}_i,v^{(\ell)}_j\rangle\bigr)^2$ is a sum of squares of
degree-$2$ polynomials, each $u^{(k,\ell)}_{i,j}$ is a real polynomial of degree exactly $4$.

We denote the full list of constraint polynomials by $p_1,\dots,p_M$, and we record their
degrees. Partition the index set into
\[
A_2=\{a:\deg p_a=2\}\quad(\text{the constraints \eqref{eq:norm},\eqref{eq:orth}}),
\qquad
A_4=\{a:\deg p_a=4\}\quad(\text{the constraints \eqref{eq:unbias}}).
\]
A configuration $X\in\mathbb{R}^{504}$ encodes seven MUBs in $\mathbb{C}^6$ if and only if
$p_a(X)=0$ for all $a$.

\subsection{Degree-$4$ SoS refutation}
We adopt the general (and most permissive) form of an SoS refutation of a system of
\emph{equalities}, which is what makes our impossibility result strongest.

\begin{definition}[Degree-$4$ SoS refutation]\label{def:refutation}
A \emph{degree-$4$ SoS refutation} of the system $\{p_a=0\}_{a=1}^M$ is a polynomial identity
in $\mathbb{R}[X]$,
\begin{equation}\label{eq:refutation}
-1\;=\;\sigma_0(X)\;+\;\sum_{a=1}^{M}\lambda_a(X)\,p_a(X),
\end{equation}
where:
\begin{itemize}
\item $\sigma_0$ is a sum of squares of real polynomials with $\deg\sigma_0\le 4$;
\item each $\lambda_a\in\mathbb{R}[X]$ is an \emph{arbitrary} real polynomial (no sign or SoS
restriction, because $p_a=0$ is an equality), subject only to the degree cap
$\deg(\lambda_a p_a)\le 4$.
\end{itemize}
\end{definition}

\begin{remark}[This subsumes the problem's literal refutation form]\label{rem:subsume}
The problem poses the refutation in Positivstellensatz form for the equality system, namely
$-1=\sigma_0+\sum_a \sigma_a\,p_a$ where each $\sigma_a$ is itself a sum of squares and
$\deg(\sigma_a p_a)\le4$. Definition~\ref{def:refutation} is strictly more permissive: a
sum-of-squares multiplier $\sigma_a$ is in particular an arbitrary real polynomial, so every
refutation in the problem's literal form is also a refutation in the sense of
Definition~\ref{def:refutation}. (More precisely, the standard way to handle an
\emph{equality} $p_a=0$ is to treat it as the two inequalities $p_a\ge0$ and $-p_a\ge0$,
contributing $\sigma_a^{+}p_a+\sigma_a^{-}(-p_a)=(\sigma_a^{+}-\sigma_a^{-})\,p_a$ with
$\sigma_a^{\pm}$ sums of squares; the resulting multiplier
$\lambda_a:=\sigma_a^{+}-\sigma_a^{-}$ is again an arbitrary real polynomial subject only to
the degree cap.) Hence proving that \emph{no} refutation of the more general
form~\eqref{eq:refutation} exists immediately rules out a refutation in the problem's stated
form. This is why establishing impossibility for Definition~\ref{def:refutation} yields the
strongest possible negative answer.
\end{remark}

The degree cap forces a structural restriction on the multipliers that is the crux of the
argument.

\begin{lemma}[Multiplier degree restriction]\label{lem:multdeg}
In any identity of the form \eqref{eq:refutation}:
\begin{enumerate}
\item for $a\in A_2$, $\deg\lambda_a\le 2$;
\item for $a\in A_4$, $\lambda_a$ is a real constant.
\end{enumerate}
\end{lemma}

\begin{proof}
For $a\in A_4$ we have $\deg p_a=4$, so $\deg(\lambda_a p_a)=\deg\lambda_a+4\le 4$ forces
$\deg\lambda_a\le 0$, i.e.\ $\lambda_a$ is a constant. (Here we use that $p_a$ in
\eqref{eq:unbias} is a nonzero polynomial of degree exactly $4$, so its degree is not lowered
upon multiplication.) For $a\in A_2$, $\deg p_a=2$ forces $\deg\lambda_a\le2$. The cap
$\deg\sigma_0\le 4$ is part of the definition. \qedhere
\end{proof}

By SoS--moment duality (made self-contained in Section~\ref{sec:duality}), the
\emph{nonexistence} of a refutation \eqref{eq:refutation} is established by producing a single
linear functional with three properties: it is normalized, it is nonnegative on squares, and it
annihilates every term $\lambda_a p_a$ that can appear. We construct such a functional in
Section~\ref{sec:pe} and verify these properties exactly.

\section{The Haar pseudo-expectation}\label{sec:pe}

\subsection{Definition}
Let $\mathrm{U}(6)$ carry its Haar probability measure. Sample seven \emph{independent}
Haar-random unitaries $U^{(1)},\dots,U^{(7)}\in\mathrm{U}(6)$, and let the columns of $U^{(k)}$
be the vectors of the $k$-th basis:
\begin{equation}\label{eq:sampling}
v^{(k)}_i \;:=\; U^{(k)} e_i\qquad(i=1,\dots,6),
\end{equation}
where $e_1,\dots,e_6$ is the standard basis of $\mathbb{C}^6$. Reading off real and imaginary
parts of the entries of the $v^{(k)}_i$ produces a random point $X\in\mathbb{R}^{504}$; let
$\mu$ be its law. Then $\mu$ is a genuine Borel probability measure on $\mathbb{R}^{504}$, with
all moments finite (the entries are bounded by $1$ in modulus). Define the linear functional
\begin{equation}\label{eq:Edef}
\widetilde{\mathbb{E}}[f]\;:=\;\mathbb{E}_{X\sim\mu}\bigl[f(X)\bigr]\;=\;\int_{\mathbb{R}^{504}} f\,d\mu ,
\qquad f\in\mathbb{R}[X].
\end{equation}
We call $\widetilde{\mathbb{E}}$ the \emph{Haar pseudo-expectation}. (We restrict attention to
$\deg f\le 4$ when we use it, but it is in fact defined on all polynomials.)

\subsection{Elementary properties from being a genuine expectation}

\begin{proposition}[Normalization and positivity]\label{prop:posnorm}
$\widetilde{\mathbb{E}}[1]=1$, and for every real polynomial $s$ one has
$\widetilde{\mathbb{E}}[s^2]\ge 0$. Consequently, for every sum of squares
$\sigma=\sum_t s_t^2$, $\widetilde{\mathbb{E}}[\sigma]\ge 0$.
\end{proposition}

\begin{proof}
$\widetilde{\mathbb{E}}[1]=\int 1\,d\mu=\mu(\mathbb{R}^{504})=1$ since $\mu$ is a probability
measure. For positivity, $s(X)^2\ge0$ pointwise, hence
$\widetilde{\mathbb{E}}[s^2]=\int s^2\,d\mu\ge0$. Linearity of the integral gives
$\widetilde{\mathbb{E}}[\sigma]=\sum_t\widetilde{\mathbb{E}}[s_t^2]\ge0$. \qedhere
\end{proof}

\subsection{The orthonormality constraints are satisfied pointwise}

\begin{proposition}[Degree-$2$ constraints]\label{prop:deg2}
For every $a\in A_2$ and every polynomial $\lambda\in\mathbb{R}[X]$,
\[
\widetilde{\mathbb{E}}[\lambda\,p_a]=0 .
\]
\end{proposition}

\begin{proof}
By construction \eqref{eq:sampling}, for each fixed $k$ the matrix $U^{(k)}$ is unitary, so its
columns $v^{(k)}_1,\dots,v^{(k)}_6$ form an orthonormal basis of $\mathbb{C}^6$ for
\emph{every} outcome in the support of $\mu$. Thus, for every $X$ in the support of $\mu$,
\[
n^{(k)}_i(X)=0,\qquad o^{(k),\mathrm{Re}}_{i,j}(X)=0,\qquad o^{(k),\mathrm{Im}}_{i,j}(X)=0
\]
for all admissible $k,i,j$; that is, $p_a(X)=0$ for all $a\in A_2$ on $\operatorname{supp}\mu$.
Hence $\lambda(X)\,p_a(X)=0$ for $\mu$-almost every $X$, and so
$\widetilde{\mathbb{E}}[\lambda p_a]=\int \lambda\,p_a\,d\mu=0$. \qedhere
\end{proof}

\subsection{The unbiasedness constraints are satisfied in expectation}

The key quantitative input is the second moment of an inner product against a Haar column.

\begin{lemma}[Uniform sphere second moment]\label{lem:sphere}
Let $w\in\mathbb{C}^6$ be the first column of a Haar-random unitary $U\in\mathrm{U}(6)$ (so $w$
is uniformly distributed on the unit sphere of $\mathbb{C}^6$). Then
\[
\mathbb{E}\bigl[\,w\,w^{*}\,\bigr]=\tfrac16\,I_6 ,
\]
and consequently, for every fixed (deterministic) unit vector $u\in\mathbb{C}^6$,
\[
\mathbb{E}\bigl[\,|\langle u,w\rangle|^2\,\bigr]=\tfrac16 .
\]
\end{lemma}

\begin{proof}
Set $S:=\mathbb{E}[ww^{*}]$, a $6\times6$ Hermitian positive semidefinite matrix. By the right
invariance of Haar measure, for any fixed $V\in\mathrm{U}(6)$ the law of $VU$ equals the law of
$U$, hence the law of $Vw$ (the first column of $VU$) equals the law of $w$. Therefore
\[
V\,S\,V^{*}=\mathbb{E}\bigl[(Vw)(Vw)^{*}\bigr]=\mathbb{E}[ww^{*}]=S
\qquad\text{for all }V\in\mathrm{U}(6).
\]
A matrix commuting with all unitaries is a scalar multiple of the identity (Schur's lemma for
the defining representation of $\mathrm{U}(6)$, which is irreducible over $\mathbb{C}$), so
$S=c\,I_6$. Taking the trace, $\operatorname{tr}S=\mathbb{E}[\operatorname{tr}(ww^{*})]
=\mathbb{E}[\,\|w\|^2\,]=1$ because $w$ is a unit vector; also $\operatorname{tr}(cI_6)=6c$.
Hence $c=\tfrac16$ and $S=\tfrac16 I_6$. Finally, for a fixed unit vector $u$,
\[
\mathbb{E}\bigl[|\langle u,w\rangle|^2\bigr]
=\mathbb{E}\bigl[u^{*}w\,w^{*}u\bigr]
=u^{*}\,S\,u=\tfrac16\,u^{*}u=\tfrac16. \qedhere
\]
\end{proof}

\begin{proposition}[Degree-$4$ constraints]\label{prop:deg4}
For every $a\in A_4$ (an unbiasedness constraint \eqref{eq:unbias}) one has
\[
\widetilde{\mathbb{E}}[p_a]=0 .
\]
\end{proposition}

\begin{proof}
Fix $k\ne\ell$ and indices $i,j$, and consider $p_a=u^{(k,\ell)}_{i,j}$. Under $\mu$, the
unitaries $U^{(k)}$ and $U^{(\ell)}$ are independent. Condition on $U^{(k)}$; then
$v^{(k)}_i=U^{(k)}e_i$ is a fixed unit vector $u$, while $v^{(\ell)}_j=U^{(\ell)}e_j$ is the
$j$-th column of an independent Haar unitary, hence uniformly distributed on the unit sphere of
$\mathbb{C}^6$ (by the left invariance of Haar measure, every column of a Haar unitary is
uniform on the sphere). Applying Lemma~\ref{lem:sphere} with $w=v^{(\ell)}_j$,
\[
\mathbb{E}\Bigl[\,\bigl|\langle v^{(k)}_i,v^{(\ell)}_j\rangle\bigr|^2 \,\Big|\, U^{(k)}\Bigr]
=\tfrac16 .
\]
The right-hand side is the deterministic constant $\tfrac16$, so taking the (outer)
expectation over $U^{(k)}$ as well gives
$\mathbb{E}\bigl[|\langle v^{(k)}_i,v^{(\ell)}_j\rangle|^2\bigr]=\tfrac16$. Therefore
\[
\widetilde{\mathbb{E}}[p_a]
=\mathbb{E}\Bigl[\,\bigl|\langle v^{(k)}_i,v^{(\ell)}_j\rangle\bigr|^2-\tfrac16\,\Bigr]
=\tfrac16-\tfrac16=0 . \qedhere
\]
\end{proof}

\section{From the pseudo-expectation to nonexistence}\label{sec:duality}

We now show directly, with no further machinery, that the existence of the Haar
pseudo-expectation rules out any degree-$4$ refutation. This is the elementary direction of
SoS--moment duality, and we give the full argument so the proof is self-contained.

\begin{theorem}[Main result]\label{thm:main}
There is no degree-$4$ SoS refutation of the MUB system \eqref{eq:norm}--\eqref{eq:unbias} in
the sense of Definition~\ref{def:refutation}. Equivalently, the answer to the problem's
question is \emph{no}.
\end{theorem}

\begin{proof}
Suppose, for contradiction, that a degree-$4$ SoS refutation \eqref{eq:refutation} exists:
\[
-1\;=\;\sigma_0(X)\;+\;\sum_{a=1}^{M}\lambda_a(X)\,p_a(X),
\]
with $\sigma_0$ a sum of squares of degree at most $4$ and $\deg(\lambda_a p_a)\le4$ for every
$a$. This is an identity of polynomials, hence it holds after applying the linear functional
$\widetilde{\mathbb{E}}$ to both sides:
\begin{equation}\label{eq:apply}
\widetilde{\mathbb{E}}[-1]
=\widetilde{\mathbb{E}}[\sigma_0]
+\sum_{a=1}^{M}\widetilde{\mathbb{E}}[\lambda_a\,p_a].
\end{equation}
We evaluate each piece.

\smallskip
\noindent\emph{Left-hand side.} By linearity and Proposition~\ref{prop:posnorm},
$\widetilde{\mathbb{E}}[-1]=-\widetilde{\mathbb{E}}[1]=-1$.

\smallskip
\noindent\emph{The square term.} $\sigma_0$ is a sum of squares, so by
Proposition~\ref{prop:posnorm},
\[
\widetilde{\mathbb{E}}[\sigma_0]\ge 0 .
\]

\smallskip
\noindent\emph{The constraint terms.} Fix $a$.
\begin{itemize}
\item If $a\in A_2$, then by Proposition~\ref{prop:deg2} (which holds for \emph{any}
multiplier $\lambda_a$, of any degree and any sign), $\widetilde{\mathbb{E}}[\lambda_a p_a]=0$.
\item If $a\in A_4$, then by Lemma~\ref{lem:multdeg}(2) the multiplier $\lambda_a$ is a real
constant $c_a$. Hence, using linearity and Proposition~\ref{prop:deg4},
\[
\widetilde{\mathbb{E}}[\lambda_a p_a]=c_a\,\widetilde{\mathbb{E}}[p_a]=c_a\cdot 0=0 .
\]
(The sign of $c_a$ is irrelevant precisely because $\widetilde{\mathbb{E}}[p_a]=0$ exactly.)
\end{itemize}
Thus $\sum_{a}\widetilde{\mathbb{E}}[\lambda_a p_a]=0$.

\smallskip
Substituting these three evaluations into \eqref{eq:apply} yields
\[
-1\;=\;\widetilde{\mathbb{E}}[\sigma_0]\;+\;0\;\ge\;0 ,
\]
i.e.\ $-1\ge0$, a contradiction. Therefore no identity of the form \eqref{eq:refutation}
exists, which is exactly the assertion that there is no degree-$4$ SoS refutation. \qedhere
\end{proof}

\section{Remarks}\label{sec:remarks}

\paragraph{Interpretation as a pseudo-expectation.}
The functional $\widetilde{\mathbb{E}}$ of \eqref{eq:Edef} is, in the language of the SoS
hierarchy, a valid \emph{degree-$4$ pseudo-expectation} for the MUB system: it is normalized
(Proposition~\ref{prop:posnorm}), its degree-$4$ moment matrix is positive semidefinite
because $\widetilde{\mathbb{E}}[s^2]\ge0$ for every polynomial $s$ with $\deg s\le2$
(Proposition~\ref{prop:posnorm}), and it satisfies the constraints to degree $4$ in the precise
sense required: $\widetilde{\mathbb{E}}[q\,p_a]=0$ whenever $\deg(q\,p_a)\le4$. Indeed for
$a\in A_2$ this holds for all $q$ with $\deg q\le2$ by Proposition~\ref{prop:deg2}, and for
$a\in A_4$ it holds for $q$ a constant by Proposition~\ref{prop:deg4}; these are exactly the
products of total degree at most $4$. The proof of Theorem~\ref{thm:main} is the elementary
``weak duality'' showing that any such functional obstructs a refutation.

\paragraph{Why the obstruction succeeds.}
The mechanism is that the orthonormality constraints \eqref{eq:norm}--\eqref{eq:orth} can be
satisfied \emph{exactly and pointwise} by a genuine random configuration (seven independent
orthonormal bases), while the only remaining constraints, the unbiasedness equalities
\eqref{eq:unbias}, are degree-$4$ and therefore, by Lemma~\ref{lem:multdeg}, may be multiplied
only by \emph{constants} in a degree-$4$ refutation. A constant multiplier sees only the global
average $\widetilde{\mathbb{E}}[p_a]$, and for independent Haar bases that average is exactly
zero by Lemma~\ref{lem:sphere}. Thus the degree-$4$ relaxation cannot detect the local
obstruction (whether the average $1/6$ can be achieved \emph{simultaneously and exactly} for
all pairs), which is the genuinely hard content of the MUB existence problem in dimension $6$.
This is consistent with the fact that the existence of seven MUBs in $\mathbb{C}^6$ remains a
long-standing open problem: a low-degree SoS refutation would have settled it.

\paragraph{Robustness of the conclusion.}
The argument used only: (a) $\mu$ is a probability measure; (b) its support consists of genuine
orthonormal bases; (c) $\widetilde{\mathbb{E}}[\,|\langle v^{(k)}_i,v^{(\ell)}_j\rangle|^2\,]=1/6$
for all $k\ne\ell,i,j$. Any probability measure with these three properties (for instance, one
need not even use Haar measure on each basis, only a measure each of whose marginals on a basis
is a genuine ONB and whose distinct bases are ``isotropically unbiased on average'') yields the
same conclusion. Haar measure is simply the most transparent choice for which
property (c) is exact.

\paragraph{On the ``natural reformulation'' clause.}
The problem allows the constraints to be replaced by ``their natural reformulation''.
Two points make our conclusion robust to this. First, the unbiasedness relation
$|\langle v^{(k)}_i,v^{(\ell)}_j\rangle|^2=\tfrac16$ is genuinely a degree-$4$ condition in
the real variables $X$: the complex inner product $\langle v^{(k)}_i,v^{(\ell)}_j\rangle$ is
itself not a real polynomial, and its squared modulus is the sum of squares of the two
degree-$2$ polynomials $\operatorname{Re}\langle\cdot\rangle$ and
$\operatorname{Im}\langle\cdot\rangle$; there is no \emph{real} polynomial \emph{equation} of
degree below $4$ in $X$ whose zero set on the orthonormality variety equals the unbiasedness
locus, so any faithful encoding without auxiliary variables uses degree-$4$ equalities, exactly
as in \eqref{eq:unbias}. Second, and decisively, our functional $\widetilde{\mathbb{E}}$ is the
expectation over a measure $\mu$ supported on genuine seven-tuples of orthonormal bases for
which, in addition, $\widetilde{\mathbb{E}}[\,|\langle v^{(k)}_i,v^{(\ell)}_j\rangle|^2\,]=1/6$
for all $k\ne\ell,i,j$. Hence $\widetilde{\mathbb{E}}$ annihilates (to the appropriate degree)
\emph{any} system of polynomial constraints whose common zero set is the MUB variety and whose
orthonormality part is degree $2$ while the unbiasedness part is degree $4$: the orthonormality
constraints vanish pointwise on $\operatorname{supp}\mu$ (so any multiplier kills them), and the
degree-$4$ unbiasedness constraints, capped to constant multipliers, are killed because their
$\widetilde{\mathbb{E}}$-average is $0$. Thus the negative answer is insensitive to the precise
algebraic presentation of the constraints, provided it respects these degrees.

\end{document}
